\scanpage{82}
\mspage{26}

qu'il convient de prendre. (Dans $\underline{\mathrm{Fib}}_0\,S$ et $\underline{\mathrm{Cofib}}_{\unsure{0}}\,S$, \inserted{ou $\underline{\mathrm{Fibsc}}_0(S)$ et $\mathrm{Cofibsc}_0(S)$,} les deux équivalences sont identiques, puisque ces \unsure{quatre} catégories sont contenues dans $\underline{\mathrm{Fib\,hot}}\,S$.) La suite nous montrera s'il en est bien ainsi.

Par contre, dans $\underline{\mathrm{Fibsc}}\,S$ et $\underline{\mathrm{Cofibsc}}\,S$, il n'y a pas de doute, suivant les principes généraux de construction d'un dérivateur à partir d'une $(\unsure{\mathcal M|M},W)$, qu'il faut prendre l'équivalence fine, fibre par fibre. À ce sujet, il faut faire attention que les foncteurs verticaux du diagramme p.~8
\[
\underline{\mathrm{Fibsc}}\,S\mathbin{\unsure{-|\longrightarrow}}\underline{\mathrm{Fib}}\,S,\qquad
\underline{\mathrm{Cofibsc}}\,S\mathbin{-}\underline{\mathrm{Cofib}}\,S,
\]
quoique \deleted{\unsure{ont}} ess. surjectifs \inserted{et fidèles}, ne sont pas pl. fidèles. Ils le seraient, \inserted{\unsure{et =|=}\ des équiv. de cat.} si dans $\underline{\mathrm{Fib}}\,S$ \deleted{\unsure{et|et Cofib}} $\underline{\mathrm{Cofib}}\,S$ je me bornais aux $S$-foncteurs \uline{cartésiens} resp. \uline{cocartésiens}.

Donc on peut \unsure{interpréter}, à équivalence près, $\underline{\mathrm{Fib}}\,\underline{\mathrm{sc}}\,S$ et $\underline{\mathrm{Cofib}}\,\underline{\mathrm{sc}}\,S$ comme $\underline{\mathrm{Fib}}\,S$ et $\underline{\mathrm{Cofib}}\,S$, où on aurait \unsure{restreint} les morphismes aux seuls $S$-\uline{f}\unsure{oncteurs} \uline{cartésiens}, resp. cocartésiens. Il semblerait donc qu'en restreignant ainsi les flèches, mais en restreignant de façon plus \unsure{draconienne} encore les localiseurs, on trouve des catégories \inserted{localisées} \unsure{coïncidentes|canoniquement|équivalentes}

% archive sheet number: 81
